% model_0806.tex
% 2026-08-06：依据 model_0805 的实际结果和 model_0806 的定向检验重写。
% 阅读顺序：固定基础 -> 已检验机制 -> 失败边界 -> 当前结论。

\subsubsection{0806 当前模型：机制与检验结论}

\paragraph{先说结论。}
当前结果可以压缩成一句话：\textbf{有限工作空间能够跟随学习轨迹，但它相对完整空间的概率预测优势取决于评价协议；现有数据也不支持“越惊讶或越不确定，下一试次就替换越多规则”这两个具体动态机制。}因此，静态 FA2 继续作为有价值但尚未确认的竞争解释；FA3-M-S 和 FA3-M-U 都不升级，当前动态替换率路线停止。动态搜索范围、背景证据和离散承诺也不开放来补救这次失败。

0805 的静态有限工作空间 FA2 仍作为统一比较起点。这里的“固定”只表示模型定义不再随结果改变，并不表示有限工作空间已经被真实数据证明。0805 的原始逐试次更新评价偏向 FS-H0，而共同的自然块冻结评价同时偏向 FA2 的 NLL 和 accuracy curve。这个排序反转说明不能脱离评价协议宣布哪一个模型已经胜出。


\subsubsection{一、先把四种状态分开}

\paragraph{固定的基础。}
知觉输入、候选规则、有限活跃工作空间、规则替换过程、双通道记忆、选择读出、试次先后顺序和潜在路径边际化均沿用 0805。它们构成 FA2 与所有后续模型共同使用的基础，不在下一轮中重新调整。

\paragraph{已经得到的正面信息。}
静态 FA2 在多个平滑窗口下都比 FS-H0 更能跟随学习轨迹，FA2 产生的规则不确定性也能够稳定增加下一试次选择的留出预测。这些结果说明“有限工作空间的状态动态”和“不确定性”都值得保留，但不能单独证明有限工作空间已经胜出，也不能证明不确定性通过替换率 $m_t$ 起作用。

\paragraph{已经检验、当前停止的机制。}
FA3-M-S 让反馈惊讶度控制 $m_t$，FA3-M-U 让规则不确定性控制 $m_t$，FA2-R 允许整个活跃集合重新生成。三者都呈现了某些局部优点，但没有同时通过概率预测、正确率水平、轨迹形状和直接机制证据，因此保留为已检验的竞争模型，不继续扩展。

\paragraph{当前保持关闭的机制。}
动态 $g_t$、背景证据 CACHE/RAW、离散承诺（commitment），以及反馈惊讶度与规则不确定性的联合动态都不进入当前模型。以后只有出现直接相关的新证据时，才把其中一个作为新的竞争解释单独检验。


\subsubsection{二、已经固定的基础模型}

\paragraph{规则怎样预测类别。}
令 $s,t,k$ 表示被试、试次和条件，$\mathcal H_k$ 为冻结的候选规则集合，$\mathcal C_k$ 为类别集合。规则 $h$ 对当前刺激的类别预测为
\begin{equation}
q_{s,t,h}(c)
=P\!\left(\widetilde{\mathbf x}_{s,t}\in R_{k,h,c}
\mid \mathbf x_{s,t},g_s^{\mathrm{percept}}\right),
\qquad \sum_{c\in\mathcal C_k}q_{s,t,h}(c)=1.
\label{eq:model0806-rule-prediction}
\end{equation}
被试特异的知觉参数和由此得到的 $q$ 在模型比较前冻结，不能再用同一段选择数据调节知觉层。

\paragraph{有限活跃工作空间。}
模型假设被试当前重点考虑容量为 $M$ 的规则集合
\begin{equation}
A_{s,t}=\{a_{s,t,1},\ldots,a_{s,t,M}\},
\qquad |A_{s,t}|=M,
\end{equation}
并为每条活跃规则保存反馈前权重 $\omega^-_{s,t}(h)$、近期记忆 $F_{s,t}(h)$ 和累积记忆 $S_{s,t}(h)$。当前统一比较中的有限模型采用硬有限工作空间（HFW）：规则退出后不继续保存在线规则特异记忆，重新进入时按新规则初始化。$M$ 在 $\{3,5,7\}$ 上比较或平均，不把某一个最优整数解释为精确的人类工作记忆容量。

\paragraph{一个试次的顺序。}
模型严格按照
\begin{equation}
\begin{aligned}
&\text{截至 }t-1\text{ 的状态和反馈}
\longrightarrow (m_{s,t},g_{s,t})
\longrightarrow \text{规则集合转移}\\
&\longrightarrow \text{选择}
\longrightarrow \text{反应时（RT）}
\longrightarrow \text{可选口头报告}
\longrightarrow \text{当前反馈}\\
&\longrightarrow \text{规则证据更新}
\longrightarrow \text{供 }t+1\text{ 使用的动态输入}
\end{aligned}
\label{eq:model0806-trial-order}
\end{equation}
运行。当前反馈只能影响下一试次及以后，不能反过来改变当前选择或当前 RT 的预测。

\paragraph{规则替换分三步。}
第一步决定换多少。为避免新规则与当前规则重复，令
\begin{equation}
\widetilde M_{s,t}=\min\{M,|\mathcal H_k|-M\},
\qquad
K_{s,t}\sim\operatorname{Binomial}(\widetilde M_{s,t},m_{s,t}).
\label{eq:model0806-replacement-count}
\end{equation}
$m_{s,t}$ 控制预期替换比例，但它不等于外部观察到的选择切换率。

第二步决定哪些规则退出。槽位 $i$ 的退出概率满足
\begin{equation}
R^{\mathrm{out}}_{s,t}(i)
\propto 1-\omega^+_{s,t-1}(a_{s,t-1,i})+\epsilon_{\mathrm{out}}.
\label{eq:model0806-exit-kernel}
\end{equation}
因此证据较弱的规则更容易退出，但退出过程仍然是随机的。

第三步决定新规则从哪里来。$K_{L,k}(h'\mid h)$ 是由冻结规则距离构造的局部搜索核，$p_{0,k}(h')$ 是全局规则先验。对当前非活跃规则定义
\begin{align}
Q_{L,s,t}(h')
&\propto I(h'\notin A_{s,t-1})
\sum_{h\in A_{s,t-1}}\omega^+_{s,t-1}(h)K_{L,k}(h'\mid h),\\
Q_{G,s,t}(h')
&\propto I(h'\notin A_{s,t-1})p_{0,k}(h'),\\
Q_{s,t}(h')
&=(1-g_{s,t})Q_{L,s,t}(h')+g_{s,t}Q_{G,s,t}(h').
\label{eq:model0806-newcomer-proposal}
\end{align}
$g_{s,t}$ 控制局部搜索和全局搜索的比例，也就是“换得多远”。在当前 FA2 和 FA3-M-U 比较中固定 $g_{s,t}=0.35$。

\paragraph{新规则怎样继承状态。}
每条新规则进入一个退出槽位并继承该槽位的概率质量。保留规则继续携带原来的近期--累积记忆差异；新规则的这一差异初始化为零。令
\begin{equation}
\Delta_{s,t-1}(h)=S_{s,t-1}(h)-F_{s,t-1}(h),
\end{equation}
保留规则使用 $\Delta^*_{s,t-1}(h)=\Delta_{s,t-1}(h)$，新规则使用 $\Delta^*_{s,t-1}(h)=0$。转移后的两条记忆同步为
\begin{align}
F_{s,t^-}(h)
&=\log\omega^-_{s,t}(h)-w_{0,s,k}\Delta^*_{s,t-1}(h),\\
S_{s,t^-}(h)
&=\log\omega^-_{s,t}(h)+(1-w_{0,s,k})\Delta^*_{s,t-1}(h).
\label{eq:model0806-memory-sync}
\end{align}
这样，规则转移和记忆模块使用同一个反馈前权重，不会在转移时凭空增加或丢失概率质量。

\paragraph{反馈怎样更新规则。}
令 $\mathcal Y_k(y_{s,t},r_{s,t})$ 表示当前选择和反馈仍允许的真实类别。规则 $h$ 得到的反馈支持为
\begin{equation}
L_{s,t}(h)
=\sum_{c\in\mathcal Y_k(y_{s,t},r_{s,t})}q_{s,t,h}(c),
\qquad \widetilde L_{s,t}(h)=\max\{L_{s,t}(h),\epsilon\}.
\label{eq:model0806-feedback-support}
\end{equation}
反馈到达后，只更新当前活跃规则：
\begin{align}
F_{s,t}(h)
&=\gamma_{s,k}F_{s,t^-}(h)+\log\widetilde L_{s,t}(h),\\
S_{s,t}(h)
&=S_{s,t^-}(h)+\log\widetilde L_{s,t}(h),\\
\ell_{s,t}(h)
&=w_{0,s,k}S_{s,t}(h)+(1-w_{0,s,k})F_{s,t}(h),\\
\omega^+_{s,t}(h)
&=\frac{\exp[\ell_{s,t}(h)]}
{\sum_{u\in A_{s,t}}\exp[\ell_{s,t}(u)]}.
\label{eq:model0806-active-update}
\end{align}
$F$ 允许旧证据逐渐衰减，$S$ 保存较稳定的累积证据。两条通道的组合方式在直接父子模型之间保持相同。

\paragraph{选择怎样产生。}
给定潜在路径或粒子 $r$，当前活跃规则先产生核心类别概率
\begin{equation}
p^{\mathrm{core},(r)}_{s,t}(c)
=\sum_{h\in A^{(r)}_{s,t}}
\omega^{-,(r)}_{s,t}(h)q_{s,t,h}(c).
\end{equation}
决策敏感度 $\kappa$ 和静态随机选择污染 $\lambda$ 再给出观测概率
\begin{align}
p^{\mathrm{read},(r)}_{s,t}(c)
&=\frac{[p^{\mathrm{core},(r)}_{s,t}(c)]^{\kappa_{s,k}}}
{\sum_{v\in\mathcal C_k}[p^{\mathrm{core},(r)}_{s,t}(v)]^{\kappa_{s,k}}},\\
p^{\mathrm{obs},(r)}_{s,t}(c)
&=(1-\lambda_{s,k})p^{\mathrm{read},(r)}_{s,t}(c)
+\frac{\lambda_{s,k}}{|\mathcal C_k|}.
\label{eq:model0806-choice-readout}
\end{align}
$\kappa$ 和 $\lambda$ 在 FA2 与所有动态子模型中使用相同的支持和边际化规则。不能只给复杂模型更大的随机污染来改善拟合，也不能把较大的 $\lambda$ 直接解释为稳定的注意特质。

\paragraph{不可见路径必须被平均。}
活跃规则集合的具体转移路径不可观察，因此选择似然必须对全部粒子路径作边际化，不能只保留事后最符合真实选择的一条路径。模型比较使用相同的粒子数、随机种子、参数支持和干扰参数处理，使父子模型之间只相差被检验的机制。


\subsubsection{三、FA3-M-U：门控通过，但完整机制未通过}

\paragraph{它要解释什么。}
FA2 使用恒定替换率
\begin{equation}
m_{s,t}=m_{s,k},
\qquad g_{s,t}=0.35,
\qquad \rho=0.
\label{eq:model0806-fa2-static}
\end{equation}
这意味着被试在学习早期和晚期都以相同比例替换规则。0805 的有限模型能跟随部分局部变化，却普遍低估学习后期正确率。一个可能原因是：静态 $m$ 在被试已经较确定时仍持续扰动正确规则。

\paragraph{规则不确定性是什么。}
反馈后规则不确定性定义为
\begin{equation}
U^{\mathrm{rule}}_{s,t}
=-\frac{\sum_{h\in A_{s,t}}
\omega^+_{s,t}(h)\log\omega^+_{s,t}(h)}{\log M}.
\label{eq:model0806-rule-uncertainty}
\end{equation}
它接近 1 时，表示多条规则仍然同样可信；它接近 0 时，表示信念已经集中在少数规则上。该量只在当前反馈以后形成，因此只能影响下一试次。

\paragraph{FA3-M-U 怎样改变替换率。}
令 $a^m_{s,t}=\operatorname{logit}(m_{s,t})$。当前候选采用
\begin{align}
a^m_{s,t+1}
=\;&\mu_{m,k}+u_{m,s}
+\phi_m\!\left[a^m_{s,t}-(\mu_{m,k}+u_{m,s})\right]\nonumber\\
&+\beta_{mU,k}\widetilde U^{\mathrm{rule}}_{s,t},
\qquad |\phi_m|<1,
\qquad \beta_{mU,k}\geq0.
\label{eq:model0806-dynamic-m-u}
\end{align}
其中 $\widetilde U^{\mathrm{rule}}$ 使用训练段冻结的中心和尺度进行标准化。通俗地说：反馈后越拿不准，下一试次越愿意换规则；越确定，下一试次越少换规则。若 $\beta_{mU,k}=0$ 且初始状态位于个体均值，该模型精确退化为静态 FA2。

\paragraph{这一轮故意不改什么。}
检验 FA3-M-U 时，$g=0.35$、$\rho=0$、$M$ 的支持、规则几何、双通道记忆、选择污染和粒子边际化全部保持不变。这样才能判断预测变化是否确实来自规则不确定性，而不是来自同时增加的其他自由度。

\paragraph{第一步：低成本门控通过。}
先冻结静态 FA2 的状态，再检验上一试次的不确定性能否增加下一试次选择的滚动留出预测。模型同时控制了 FA2 的正确选择概率、练习、刺激歧义性和上一试次对错。31 名被试的被试等权 $\Delta\mathrm{LPD}$/试次为 $0.00344$，bootstrap 95\% 区间为 $[0.00033,0.00711]$，$22/31$ 名被试改善，$11/11$ 个预测折叠的系数为正。更换独立粒子种子后，增量为 $0.00381$，区间仍完全高于零。因此，不确定性含有稳定的额外选择信息。

\paragraph{第二步：合成恢复允许进入真实数据。}
在 16 份合成数据上，仅用 choice 时，静态--动态家族恢复率为 $75\%$；加入与同一潜在路径联合生成的理想 RT 后，恢复率为 $100\%$。有效候选数平均从 $5.34$ 降至 $2.26$，减少 $57.7\%$。在动态真值数据中，动态效应方向的恢复率为 $100\%$。这说明该检验流程在理想条件下能够发现 FA3-M-U，也再次说明 RT 原则上可以明显缩小等价集合。

\paragraph{第三步：真实 choice 没有通过完整标准。}
真实数据的滚动比较包含 31 名被试和 125 个预测块。隔离重跑实际生成了预期的 16,128 个模型分量，没有缺失或多余文件，也只包含预先指定的两组粒子种子。把两组种子一起边际化后，FA3-M-U 相对静态 FA2 的被试等权 $\Delta\mathrm{NLL}$/试次为 $0.00332$，bootstrap 95\% 区间为 $[-0.00068,0.00707]$，$21/31$ 名被试改善。分开计算两组种子时，增量分别为 $0.00297$ 和 $0.00350$，95\% 区间分别为 $[-0.00146,0.00712]$ 和 $[-0.00050,0.00729]$。两组区间都包含零，因此这个小幅改善还不能排除随机波动。

更重要的是，两组种子都显示局部轨迹变差：16 试次局部曲线相关分别从 $0.489$ 降到 $0.380$，以及从 $0.481$ 降到 $0.388$。合并两组种子后，相关从 $0.488$ 降到 $0.383$。FA3-M-U 确实稍微减轻了晚期正确率低估，早期、中期和晚期的平均 NLL 也都略有改善；但它没有可靠改善总体预测，并且稳定地损害了局部学习轨迹的跟随。根据预先规定的完整标准，这一项仍然判为不通过。

统一多指标复核得到同样的结果。相对使用相同参数网格的静态 FA2，FA3-M-U 的 NLL/试次平均降低 $0.00246$，但 95\% 区间为 $[-0.00640,0.00175]$，不能排除随机波动。曲线 MAE 和 centered MAE 都没有可靠改善，曲线相关则平均降低 $0.0795$，95\% 区间为 $[-0.1427,-0.0315]$。换句话说，它让预测曲线动得更多，却没有更准确地跟随真实学习曲线。

\paragraph{第四步：真实 RT 不支持预期的搜索成本。}
若 FA3-M-U 的解释正确，预测替换量越大，RT 应越长。在与上述隔离重跑对应的 RT 检验中，加入预测替换量后的 $\Delta\mathrm{LPD}$/试次为 $0.00258$，bootstrap 95\% 区间为 $[-0.00040,0.00625]$。替换量系数的折叠中位数为 $-0.0098$，只有 $5/11$ 个折叠为正；只保留至少 5 名被试的稳定折叠后，中位数为 $-0.0743$，只有 $1/7$ 个折叠为正。因此，真实 RT 没有支持“替换更多规则需要更多时间”这一解释。

\paragraph{综合判定。}
低成本门控只证明“不确定性有预测信息”；完整 FA3-M-U 检验才回答“这些信息是否通过提高替换率发挥作用”。后一个问题的答案是否定的。因此 FA3-M-U 不升级，当前动态 $m_t$ 路线停止，也不再运行自主生成来为已经失败的机制寻找补救解释。


\subsubsection{四、已经检验、当前停止的机制}

\paragraph{0805 给出的基准事实。}
0805 的正式运行覆盖实验条件 1 的仅选择模型。在最终时间留出段，完整空间基线 FS-H0 的被试等权平均 NLL/试次为 $0.2295$，有限模型中最好的 FA2-$M=3$ 为 $0.2567$，只有 $3/32$ 名被试的 FA2 优于 FS-H0；NLL 越低表示预测越好。因此，现有选择数据没有支持有限工作空间在总体预测上取代完整空间基线。

有限模型仍捕捉到另一部分结构：在 16 试次窗口下，外部留出段的平均正确率曲线相关从 FS-H0 的 $0.187$ 提高到 FA2-$M=3$ 的 $0.499$。但是，FA2 普遍低估学习后期正确率。FS-H0 则给出接近 $0.90$ 的较平坦预测，并且 $28/32$ 名被试的最终 $\lambda$ 到达预设上界 $0.20$。因此两个模型各自遗漏了不同的时间结构，不能只根据其中一个指标宣布有限工作空间已经成立。

\paragraph{统一多指标复核表明，模型排序依赖评价协议。}
新复核把 FS-H0、FA2、FA2-R、FA3-M-S 和 FA3-M-U 放在相同的自然块滚动评价中：每个 64 试次评价块开始时冻结参数分量权重，评价块内不再利用新选择调整这些权重。31 名被试共提供 8000 个共同评价试次。在这一更严格的块间迁移口径下，FS-H0 与 FA2-$M=3$ 的 NLL/试次分别为 $0.4772$ 和 $0.4104$，16 试次 accuracy curve 的 MAE 分别为 $0.1244$ 和 $0.0874$，曲线相关分别为 $0.103$ 和 $0.727$。也就是说，FA2 在这个口径下同时改善了概率预测和学习轨迹。

这个结果不推翻 0805，而是说明两种评价回答了不同问题。0805 允许模型在最终留出段内逐试次更新参数分量权重，检验在线一步预测；自然块评价则在整块内冻结权重，检验模型能否把此前学到的参数迁移到下一块。FS-H0 与 FA2 的 NLL 排序在两种口径下反转，因此现有数据不足以给出不依赖评价协议的总体胜者。以后必须对所有模型使用相同协议，并同时报告逐试次更新和自然块冻结结果。

\paragraph{简单概率校准不能解释这个反转。}
复核还只改变选择概率的截距和斜率，不重新运行规则学习，也不改变任何潜在状态。校准参数只使用每个评价块以前的数据。在自然块滚动评价中，FA2-$M=3$ 的 NLL/试次反而增加 $0.01173$，95\% 区间为 $[0.00437,0.01876]$。在 0805 原始拆分中，使用 inner fit 拟合并由 inner validation 选择校准后，outer NLL 变化为 $0.00113$，区间为 $[-0.00161,0.00397]$；校准后的 FA2 相对原始 FS-H0 仍差 $0.02835$，区间为 $[0.01730,0.03948]$。因此，当前这种简单的单调概率校准不足以把 NLL 冲突归因于读出失准，但这不排除其他读出模型。

\paragraph{FA3-M-S：惊讶度动态已经检验。}
反馈惊讶度定义为
\begin{equation}
S^{\mathrm{fb}}_{s,t}
=-\log\max\left\{
\sum_{h\in A_{s,t}}\omega^-_{s,t}(h)L_{s,t}(h),\epsilon
\right\}.
\label{eq:model0806-feedback-surprise}
\end{equation}
FA3-M-S 假设反馈越出乎预料，下一试次的 $m_t$ 越高。该模型已经实现，并通过了静态端点和无反馈泄漏检查。

在 16 份自主生成数据上，仅用选择时，静态--动态家族恢复率为 $12/16=75\%$；加入与同一潜在搜索路径联合生成的理想 RT 后，恢复率为 $16/16=100\%$。有效候选数平均从 $7.13$ 降至 $3.55$。这说明设计良好的 RT 在原则上可以缩小等价集合。但在动态生成数据中，选择+RT 下仍平均保留约 $5.47$ 个有效候选，而且 RT 参数固定在生成真值，因此这是理想条件下的识别上限，不是真实 RT 已支持该机制。

\paragraph{真实选择没有支持 FA3-M-S。}
实验条件 1 的滚动时间比较包含 31 名可评价被试和 125 个预测块。静态 FA2 与动态 FA3-M-S 的被试等权 NLL/试次分别为 $0.4103$ 和 $0.4079$，平均改善只有 $0.0024$，被试层自助法（bootstrap）95\% 区间为 $[-0.0013,0.0062]$。至少有 3 个预测块的 20 名被试平均反而恶化 $0.0024$，合并全部预测试次后也恶化 $0.0012$。

FA3-M-S 确实把替换率从早期约 $0.157$ 降到晚期约 $0.070$，并把平均正确率低估从 $5.4\%$ 减少到 $4.0\%$。但是，16 试次局部曲线相关从 $0.487$ 降到 $0.438$。也就是说，它让晚期平均正确率更接近真实值，却更难跟随真实选择的局部变化。

统一多指标复核给出同一判断。相对使用相同参数网格的静态 FA2，FA3-M-S 的 NLL/试次平均改善 $0.00231$，但 95\% 区间为 $[-0.00138,0.00599]$，仍包含零。它把曲线 MAE 降低 $0.00409$，主要因为平均低估减少；去掉整体高低差后的 centered MAE 反而增加 $0.00316$，曲线相关也降低 $0.0398$。因此它改善的是曲线水平，而不是局部起伏的时间形状。

\paragraph{惊讶度与下一选择的直接连接也未通过。}
在控制静态 FA2 的正确率预测、练习、刺激歧义性和上一试次对错以后，上一试次惊讶度带来的被试等权 $\Delta\mathrm{LPD}$/试次为 $0.00010$，95\% 区间为 $[-0.00048,0.00087]$，只有 $13/31$ 名被试改善。标准化惊讶度系数在第 2--3 块平均为 $-0.0766$，第 4 块以后才变为 $0.0600$。因此现有数据不支持一个跨学习阶段保持不变的正向惊讶度效应。

\paragraph{RT 只有弱方向信号。}
使用完全由选择冻结的状态后，预测替换量的 RT 系数在 $11/11$ 个折叠中为正，方向上符合“换得越多，反应越慢”。但其留出 $\Delta\mathrm{LPD}$/试次只有 $0.00054$，95\% 区间为 $[-0.00268,0.00383]$；正的 $\Delta\mathrm{LPD}$ 表示新增变量改善预测。允许被试特异的练习斜率后，增量为 $0.00404$，区间仍包含零。这个方向可作为新数据中的预注册预测，但不足以在当前数据上继续修改惊讶度动态方程。

\paragraph{FA2-R：整体重置也不继续扩展。}
FA2-R 允许以概率 $\rho$ 放弃整个活跃集合并重新生成。0805 中，FA2-R 在 $M=3$ 和 $M=5$ 时比对应 FA2 更差，在 $M=7$ 时只有很小差异。统一复核进一步显示，相对 FA2-$M=3$，FA2-R-$M=3$ 的 NLL/试次恶化 $0.00927$，曲线 MAE 恶化 $0.00745$；但 centered MAE 改善 $0.00303$，曲线相关提高 $0.0417$。这表示整体重置稍微改善了曲线形状，却把总体正确率压得更低。它因此保留为轨迹形状上的竞争模型，但不升级，也不继续建立动态 $\rho_t$。


\subsubsection{五、已经有定义、但当前不能进入的分支}

\paragraph{动态 $g_t$：当前不进入。}
$m_t$ 回答换多少，$g_t$ 回答换得多远。只有 FA3-M-U 通过，而且 RT 或口头报告能够区分局部新规则与全局新规则时，才允许令 $g_t$ 动态变化。例如可使用
\begin{equation}
a^g_{s,t+1}
=\mu_{g,k}+u_{g,s}
+\phi_g[a^g_{s,t}-(\mu_{g,k}+u_{g,s})]
+\beta_{gS,k}\widetilde S^{\mathrm{fb}}_{s,t}
+\beta_{gU,k}\widetilde U^{\mathrm{rule}}_{s,t},
\label{eq:model0806-dynamic-g}
\end{equation}
并令 $g_{s,t}=\operatorname{logit}^{-1}(a^g_{s,t})$。FA3-M-U 已经没有通过，所以当前不进入动态 $g_t$。只有未来数据直接显示局部搜索与全局搜索之间的变化时，才把它作为单独假设检验。

\paragraph{背景证据 CACHE/RAW：回答退出规则是否还记得。}
HFW、CACHE 和 RAW 的区别只在于规则退出以后发生什么：HFW 不保存在线规则特异记忆；CACHE 保存退出时的记忆；RAW 还允许非活跃规则继续吸收反馈。只有出现“旧规则重新进入后，模型重新学习得过慢”的稳定残差，并且 RT 或口头报告提供旧规则立即恢复的证据时，才依次比较 HFW、CACHE 和 RAW。背景证据不能因为 FA2 预测较差就直接加入，也不能与 FA3-M-U 同时首次开放。

\paragraph{离散承诺：连续稳定化的竞争解释。}
FA3-M-U 用连续降低 $m_t$ 解释学习后的稳定。只有它仍系统性低估被试长时间坚持同一规则，并且连续错误、RT 或口头报告显示明确的单规则锁定状态时，才比较离散承诺。它应当作为竞争解释单独检验，而不是预先与背景证据或动态 $g_t$ 捆绑。

\paragraph{FA3-M-SU：当前不进入。}
FA3-M-SU 让惊讶度和规则不确定性同时控制 $m_t$。由于仅惊讶度分支已经失败，现在直接同时估计两个效应只会增加相互补偿和不可识别性。因此，除非未来独立数据重新支持稳定的惊讶度效应，否则不进入该联合模型。

\paragraph{口头报告仍是可选测量。}
没有口头报告的实验继续以选择+RT 为主要识别设计。有报告时，只把它作为当前活跃规则内容的额外测量；报告缺失只删除相应观测项，不改变选择、反馈更新或动态策略方程。


\subsubsection{六、RT 与不可识别性的处理原则}

\paragraph{不要求选择数据单独找出唯一参数。}
仅靠选择，$m_t$ 的个体均值、持续性和动态效应可能相互补偿。因此当前目标不是强迫每名被试只剩一个参数点，而是恢复较稳定的量：动态效应方向、反馈后下一试次的期望替换量、早期与晚期的平均替换差异，以及正确规则处于活跃集合的概率。

\paragraph{RT 应当缩小等价集合。}
令 $\mathcal E^{\mathrm{choice}}_s$ 为选择预测在预先冻结的容差内都可接受的模型--参数集合。加入 RT 后，可接受集合定义为
\begin{equation}
\mathcal E^{\mathrm{choice+RT}}_s
=\left\{e\in\mathcal E^{\mathrm{choice}}_s:
d_{\mathrm{RT}}(e,e_{\mathrm{best}})\leq\epsilon_{\mathrm{RT}}\right\}
\subseteq\mathcal E^{\mathrm{choice}}_s.
\label{eq:model0806-equivalence}
\end{equation}
这个包含关系在数学上成立，但现实 RT 不一定提供足够信息。因此必须直接报告加入 RT 前后模型混淆和有效候选数减少了多少，不能只说“理论上 RT 有帮助”。

\paragraph{RT 先检验换多少，再检验换多远。}
第一步只检查实际替换比例 $K/M$ 是否在选择不确定性、练习、刺激和运动协变量之外增加 RT 预测。只有这一步通过，才加入新规则距离来约束 $g_t$。相应的第一版 RT 模型为
\begin{align}
\log(\mathrm{RT}_{s,t})\sim t_\nu\!\Big(&
\alpha_{s,k}+\mathbf b_{0,k}^{\mathsf T}\mathbf Z_{s,t}
+b_{U,k}U^{\mathrm{choice}}_{s,t}\nonumber\\
&+b_{K,k}K_{s,t}/M
+b_{d,k}\bar d^{\mathrm{new}}_{s,t},
\sigma_{\mathrm{RT},k}\Big).
\label{eq:model0806-rt}
\end{align}
$b_K$ 不通过时，动态 $m_t$ 最多只能被解释为选择预测结构，不能获得“搜索成本”的强心理解释；$b_d$ 不通过时，不解释动态 $g_t$。本轮 FA3-M-U 的 $b_K$ 没有通过，而且稳定折叠中的方向主要为负，因此当前数据也不支持继续用 RT 检验动态 $g_t$。


\subsubsection{七、以后怎样判断一个机制是否通过}

\paragraph{先固定两种评价方式。}
第一种是逐试次更新：每看到一个新选择，模型就更新参数分量权重，再预测下一试次。它回答模型能否进行在线一步预测。第二种是自然块冻结：进入一个新的 64 试次块后，整块都使用块开始前确定的权重。它回答模型能否把前一阶段学到的东西带到下一阶段。所有模型必须在完全相同的试次上接受这两种评价。如果两种评价给出相反排序，就把结果写成“对评价协议敏感”，不能宣布某个模型已经胜出。

\paragraph{概率预测和学习曲线必须同时检查。}
NLL 和 Brier 分数检查模型给出的选择概率是否可靠；ECE 检查模型说“有多大把握”是否可信。accuracy curve 则检查模型有没有跟上真实学习过程。曲线检查至少包括整体高低误差、去掉整体高低差后的形状误差、曲线相关，以及 8、16、32、64 试次窗口下的结果。还要分别检查学习早期、中期、晚期和错误后的恢复过程。所有不确定区间都以被试为单位重采样，不能把同一名被试的大量试次当成彼此独立的证据。

\paragraph{不再用一个总分掩盖冲突。}
以后采用以下简单规则：
\begin{enumerate}
    \item NLL 改善、但学习曲线形状变差：只能说概率预测有局部改善，机制不升级；
    \item 曲线形状改善、但 NLL 或平均正确率明显变差：保留为竞争模型，不能取代父模型；
    \item choice 改善、但 RT 或口头报告的方向与机制预期相反：最多保留为预测模型，不能作强心理解释；
    \item 结果随评价协议反转：标记为协议敏感，需要新数据确认；
    \item 只有概率、曲线水平、曲线形状和直接机制证据大体一致，而且合成恢复检验能够识别该机制时，才允许升级。
\end{enumerate}
这里没有要求 choice 单独确定唯一参数。要求的是：choice 先把候选范围压小，RT 或口头报告再提供独立约束，而且新增复杂度确实带来稳定、可重复的预测收益。


\subsubsection{八、本轮分析实际完成了什么}

\begin{enumerate}
    \item 固定了 FA2 的知觉输入、规则几何、$M$ 支持、记忆、$g=0.35$、$\rho=0$、选择干扰参数和粒子设置；
    \item 不确定性的低成本选择门控通过，并由独立粒子种子复核；
    \item FA3-M-U 已经实现，且 $\beta_{mU}=0$ 时精确退化为静态 FA2；
    \item 仅 choice 和理想 choice+RT 的合成恢复已经完成，结果允许进入真实数据检验；
    \item 真实 choice 滚动预测及其独立粒子种子复核均未通过完整标准；
    \item 真实 RT 检验及其独立粒子种子复核均未通过正搜索成本标准；
    \item 按预先规定的停止规则，未继续运行自主生成，也未开放动态 $g_t$、背景证据或离散承诺；
    \item 在 31 名被试、8000 个共同评价试次上完成了统一多指标复核，同时比较 NLL、Brier、概率校准、accuracy curve、不同窗口、学习阶段和错误后恢复；
    \item 完成了前缀数据交叉验证的单调概率校准；这种简单校准没有改善 FA2，也不能解释两种评价协议下的排序反转；
    \item 确认 FS-H0 与 FA2 的相对结论对评价协议敏感，因此当前不指定唯一主模型。
\end{enumerate}

下一步不应继续给当前模型增加机制。应当先把两种评价协议和上述多指标判定标准固定下来，然后在新的独立数据上同时比较 FS-H0、静态 FA2 和已经停止的动态模型。新实验最好同时记录 choice 和 RT；条件允许时再加入少量口头报告。另一个开放问题是“不确定性为何能预测下一选择”：它可能影响选择噪声、信念读出或其他尚未测量的过程，不一定通过改变规则替换率起作用。在出现能够区分这些解释的新证据以前，不再增加动态 $m_t$ 的变体。


\subsubsection{九、当前允许的结论}

0805 的逐试次更新评价偏向 FS-H0，统一自然块冻结评价则偏向静态 FA2。两个结果都来自正确的计算，但回答的问题不同。因此当前最准确的说法不是“有限工作空间已经成立”或“有限工作空间已经被否定”，而是：静态 FA2 确实捕捉了 FS-H0 容易遗漏的学习轨迹信息，但它相对 FS-H0 的总体概率优势还没有跨评价协议确认。

两种动态替换率都没有通过完整标准。FA3-M-S 主要修正了平均正确率偏低，却损害了局部曲线形状；FA3-M-U 没有可靠改善 NLL 或曲线误差，曲线相关和 RT 方向也不支持预期解释。FA2-R 对曲线形状有一点帮助，但使 NLL 和平均水平更差。因此当前不升级任何动态替换或整体重置机制。

现在允许保留四点结论。第一，在理想模拟中，choice+RT 能够明显缩小不可识别集合。第二，有限工作空间状态含有真实的学习轨迹信息，值得继续作为竞争解释。第三，规则不确定性能够预测下一试次选择，但现有证据不支持它通过提高替换率 $m_t$ 起作用。第四，简单的概率校准不能解决 FA2 与 FS-H0 的评价冲突。边界也必须写清楚：这些数据已经参与模型设计，最终机制判断需要新的独立数据确认。
